\paragrafos{Αφαίρεση ρητών: η απόσταση ανάμεσα σε δύο θέσεις}{Α.7.4, «Αφαίρεση ρητών αριθμών»}

\begin{anagke}
Ο μετεωρολόγος στην τηλεόραση του ξενοδοχείου λέει: «Η μέγιστη θερμοκρασία σήμερα είναι $+5$ βαθμοί και η ελάχιστη $-3$. Η διαφορά είναι $8$ βαθμοί». Ο Στέλιος κάνει τον λογαριασμό: «Πέντε μείον τρία κάνει δύο. Ο μετεωρολόγος κάνει λάθος». Η Πολυξένη τον διορθώνει: «Δεν αφαιρούμε το τρία από το πέντε. Αφαιρούμε το μείον τρία!».

Η Αγγελική σχεδιάζει τον άξονα: «Από το $-3$ ως το $+5$ πόσα βήματα υπάρχουν;». Μετράνε: τρία βήματα ως το μηδέν και πέντε ακόμη. Οκτώ. Ο Νίκος συμπληρώνει: «Η αφαίρεση είναι αναποδογύρισμα της πρόσθεσης: ψάχνω αριθμό που προστιθέμενος στο $-3$ να δώσει $+5$. Δηλαδή $(-3)+x=5$. Αφού $(-3)+8=5$, το $5-(-3)$ είναι $8$».

Ο Ηλίας το γενικεύει: «Αν η αφαίρεση ψάχνει «πόσο πρέπει να προσθέσω», μήπως $\alpha-\beta$ είναι το ίδιο με το να προσθέσω τον αντίθετο του $\beta$;». Έτσι η παρέα ανακαλύπτει ότι κάθε αφαίρεση μπορεί να γίνει πρόσθεση και ότι όλα όσα ξέρουν για την πρόσθεση λειτουργούν και για την αφαίρεση.
\end{anagke}

\begin{lexi}{αφαίρεση, διαφορά}{αφαίρεση, διαφορά}
Η \lexiword{αφαίρεση} είναι το «αφ-αιρώ», «βγάζω από κάτι». Η \lexiword{διαφορά} είναι αυτό που μένει και δείχνει «πόσο διαφέρουν» δύο αριθμοί. Ο \lexiword{μειωτέος} είναι αυτός που «μειώνεται» και ο \lexiword{αφαιρετέος} αυτός που αφαιρείται.
\end{lexi}

\begin{metafrasi}
Η διαφορά $\alpha-\beta$ είναι ο αριθμός $x$ που... & $\beta+x=\alpha$\\
Αφαίρεση σημαίνει πρόσθεση του αντίθετου. & $\alpha-\beta=\alpha+(-\beta)$\\
Η διαφορά θερμοκρασιών $+5$ και $-3$. & $(+5)-(-3)=8$\\
Αφαιρώ ολόκληρη παρένθεση. & $\alpha-(\beta+\gamma)=\alpha-\beta-\gamma$\\
\end{metafrasi}

\tmima{Η θεωρία}

\begin{orismos}{διαφορά}
\textbf{Διαφορά} $\alpha-\beta$ δύο ρητών ορίζεται ως η μοναδική λύση της εξίσωσης $\beta+x=\alpha$.
\end{orismos}

Για παράδειγμα η διαφορά $(+5)-(-3)$ είναι η λύση της $(-3)+x=5$, δηλαδή $x=8$. Από τον ορισμό προκύπτει ο κανόνας:

\begin{kanonas}{αφαίρεση ρητών}
Για να αφαιρέσω έναν αριθμό, προσθέτω τον \textbf{αντίθετό} του: $$\alpha-\beta=\alpha+(-\beta).$$ Έτσι $7-(+9)=7+(-9)=-2$ και $(-4)-(-6)=(-4)+(+6)=2$.
\end{kanonas}

Η διαφορά $\alpha-\beta$ δείχνει και τον δρόμο από το $\beta$ στο $\alpha$ στον άξονα: αν $\alpha>\beta$ είναι θετική, αλλιώς αρνητική.

\tmimasmall{Αλγεβρικό άθροισμα και απαλοιφή παρενθέσεων}

Μια παράσταση με προσθέσεις και αφαιρέσεις γράφεται ως πρόσθεση ρητών. Για παράδειγμα $12-15+8-20+3=12+(-15)+8+(-20)+3$, και μπορώ να τα αθροίσω με όποια σειρά θέλω: $(12+8+3)+(-15-20)=23-35=-12$. Τέτοιο άθροισμα ονομάζεται \textbf{αλγεβρικό άθροισμα}. Για να σβήσω παρενθέσεις:
\begin{itemize}[leftmargin=1.2em,itemsep=0.3ex,topsep=2pt]
\item Αν μπροστά από την παρένθεση υπάρχει $+$, τη σβήνω χωρίς να αλλάξω τα πρόσημα: $\alpha+(\beta-\gamma)=\alpha+\beta-\gamma$.
\item Αν μπροστά υπάρχει $-$, τη σβήνω και \textbf{αλλάζω το πρόσημο όλων} των όρων μέσα: $\alpha-(\beta-\gamma+\delta)=\alpha-\beta+\gamma-\delta$.
\end{itemize}

\begin{prosoxi}
\begin{itemize}[leftmargin=1.2em,itemsep=0.3ex,topsep=2pt]
\item Η διαφορά δεν είναι αντιμεταθετική: $5-3=2$ αλλά $3-5=-2$. Δεν ισχύει $\alpha-\beta=\beta-\alpha$.
\item Η αφαίρεση δεν «μικραίνει» πάντα: $3-(-2)=5$, μεγαλύτερο από το $3$.
\item Με το $-$ μπροστά από παρένθεση αλλάζω το πρόσημο \emph{όλων} των όρων της παρένθεσης, όχι μόνο του πρώτου: $10-(3-8+2)=10-3+8-2$.
\item Η διαφορά θερμοκρασιών $+5$ και $-3$ είναι $8$, όχι $2$. Αφαιρώ τον αριθμό μαζί με το πρόσημό του.
\end{itemize}
\end{prosoxi}

\tmima{Λυμένα παραδείγματα}

\paradeigma{Η διαφορά των θερμοκρασιών}
\tixero Μέγιστη $+5$°, ελάχιστη $-3$°.\par
\tithelo Η θερμική διακύμανση, δηλαδή η διαφορά μέγιστης και ελάχιστης.\par
\possindeo $(+5)-(-3)=(+5)+(+3)=8$. Έλεγχος στον άξονα: από το $-3$ ως το $+5$ υπάρχουν $8$ μονάδες.\par
\apantisi $8$ βαθμοί.\par
\elegxos $(-3)+8=5$.

\paradeigma{Με τον κανόνα}
\tixero Οι διαφορές $7-(+9)$, $(-4)-(-6)$, $(+3{,}5)-(-1{,}5)$, $\left(-\frac{1}{2}\right)-\left(+\frac{1}{3}\right)$.\par
\tithelo Να τις υπολογίσω.\par
\possindeo Προσθέτω τον αντίθετο: $7+(-9)=-2$. $(-4)+(+6)=2$. $3{,}5+1{,}5=5$. $-\frac{1}{2}-\frac{1}{3}=-\frac{3}{6}-\frac{2}{6}=-\frac{5}{6}$.\par
\apantisi $-2$, $2$, $5$, $-\frac{5}{6}$.\par
\elegxos Για την πρώτη: $9+(-2)=7$. Σωστό.

\paradeigma{Απαλοιφή παρενθέσεων}
\tixero Η παράσταση $10-(3-8+2)$.\par
\tithelo Να την υπολογίσω με δύο τρόπους.\par
\possindeo \emph{1ος τρόπος:} πρώτα η παρένθεση: $3-8+2=-3$, και $10-(-3)=13$. \emph{2ος τρόπος:} σβήνω την παρένθεση αλλάζοντας πρόσημα: $10-3+8-2=13$.\par
\apantisi $13$.\par
\elegxos Και οι δύο δίνουν $13$.

\paradeigma{Αλγεβρικό άθροισμα}
\tixero Η παράσταση $12-15+8-20+3$.\par
\tithelo Να την υπολογίσω.\par
\possindeo Ομαδοποιώ: θετικοί $12+8+3=23$. Αρνητικοί $-15-20=-35$. $23+(-35)=-12$.\par
\apantisi $-12$.\par
\elegxos Με τη σειρά: $12-15=-3$, $-3+8=5$, $5-20=-15$, $-15+3=-12$.

\paradeigma{Εξισώσεις με αφαίρεση}
\tixero Οι εξισώσεις $x+(-3)=5$, $x-7=-2$ και $4-x=9$.\par
\tithelo Να τις λύσω.\par
\possindeo $x=5-(-3)=8$. $x=-2+7=5$. $x=4-9=-5$ (από την $\alpha-x=\beta\Rightarrow x=\alpha-\beta$).\par
\apantisi $8$, $5$, $-5$.\par
\elegxos $8+(-3)=5$, $5-7=-2$, $4-(-5)=9$.

\tmima{Ασκήσεις}

\zoni{Το κατάλαβα}{διαφορές, παρενθέσεις, αλγεβρικά αθροίσματα}

\ex Υπολόγισε: α) $(+5)-(+3)$ \quad β) $(+3)-(+5)$ \quad γ) $(-4)-(-1)$ \quad δ) $(-1)-(-4)$ \quad ε) $0-(+6)$ \quad στ) $0-(-6)$.

\ex Μετάτρεψε σε πρόσθεση και υπολόγισε: α) $7-(+9)$ \quad β) $(-4)-(-6)$ \quad γ) $(+3{,}5)-(-1{,}5)$ \quad δ) $\left(-\frac{1}{2}\right)-\left(+\frac{1}{3}\right)$ \quad ε) $(-2{,}4)-(+0{,}6)$ \quad στ) $8-(-8)$.

\ex Λύσε: α) $x+7=3$ \quad β) $x-5=-2$ \quad γ) $4-x=9$ \quad δ) $x+(-6)=-6$ \quad ε) $-3-x=2$.

\ex Κάνε απαλοιφή παρενθέσεων και υπολόγισε: α) $10-(3-8+2)$ \quad β) $5+(-3+7)-(2-9)$ \quad γ) $-(-4+5)-(7-10)$ \quad δ) $20-(-5)-(-(3))$.

\ex Υπολόγισε τα αλγεβρικά αθροίσματα: α) $12-15+8-20+3$ \quad β) $-2{,}5+7-4{,}5-1$ \quad γ) $\frac{1}{2}-\frac{3}{4}+\frac{1}{4}-1$.

\ex Σωστό ή λάθος; α) $\alpha-\beta=\beta-\alpha$. \quad β) $\alpha-0=\alpha$. \quad γ) $0-\alpha=-\alpha$. \quad δ) $(-3)-(-3)=0$. \quad ε) $(+4)-(-2)=2$. \quad στ) Η διαφορά δύο θετικών μπορεί να είναι αρνητική.

\zoni{Το χρησιμοποιώ}{θερμοκρασίες, υψόμετρα, χρόνια}

\ex \textbf{Η διακύμανση.} Βρες τη διαφορά μέγιστης και ελάχιστης θερμοκρασίας όταν α) μέγιστη $+5$, ελάχιστη $-3$ \quad β) μέγιστη $-2$, ελάχιστη $-11$.

\ex \textbf{Το υψόμετρο.} Ο Όλυμπος είναι στα $+2\,917$ μ. και η Νεκρά Θάλασσα στα $-430$ μ. Πόση είναι η υψομετρική διαφορά;

\ex \textbf{Η ιστορία.} Η δυναστεία των Χαν ξεκίνησε το $-206$ και ο Διόφαντος έζησε γύρω στο $250$ μ.Χ. Πόσα χρόνια μεσολαβούν (αγνοώντας ότι δεν υπήρξε έτος $0$);

\ex \textbf{Το ταμείο.} Η Ζωή έχει $-15$ € και ο Ηλίας $+8$ €. Πόση είναι η διαφορά των δύο ποσών;

\ex \textbf{Ο δύτης.} Ο δύτης ανεβαίνει από το $-22$ μ. στο $-9$ μ. Πόσα μέτρα ανέβηκε;

\ex \textbf{Ο ανελκυστήρας.} Από τον όροφο $3$ στον $-2$ η μεταβολή είναι $(-2)-3$. Πόσο είναι και τι σημαίνει;

\ex \textbf{Η πτώση.} Η θερμοκρασία από $+4$° έπεσε στα $-7$°. Πόση είναι η μεταβολή (νέα $-$ παλιά);

\zoni{Το σκέφτομαι}{όταν η αφαίρεση μεγαλώνει}

\ex Ο Στέλιος γράφει $5-(-3)=2$. Βρες το λάθος και δείξε το σωστό με τον άξονα.

\ex Μικραίνει πάντα ένας αριθμός όταν του αφαιρώ κάτι; Δώσε παράδειγμα.

\ex Αν $\alpha>\beta$, ποιο πρόσημο έχει η διαφορά $\alpha-\beta$; Και αν $\alpha<\beta$;

\ex Δείξε με $\alpha=10$, $\beta=3$, $\gamma=4$ ότι $\alpha-(\beta+\gamma)=\alpha-\beta-\gamma$.

\ex Βρες τον αριθμό που, αν του αφαιρέσουμε το $-5$, δίνει $2$.

\zoni{Για το σπίτι}{διαφορές του κόσμου}

\ex Βρες τη διαφορά θερμοκρασίας ανάμεσα στο πιο ζεστό και στο πιο κρύο μέρος του κόσμου σήμερα.

\ex Φτιάξε μια παράσταση με έξι όρους και παρενθέσεις και υπολόγισέ την με δύο τρόπους.
